% !TeX program = xelatex
% UTF-8. Compile twice with XeLaTeX; all figures are drawn in this file.
\documentclass[UTF8,a4paper,11pt]{ctexart}
\usepackage[left=23mm,right=23mm,top=24mm,bottom=23mm,headheight=15pt]{geometry}
\usepackage{amsmath,amssymb,booktabs,tabularx,array,xcolor,tikz,fancyhdr}
\usetikzlibrary{arrows.meta,positioning,calc}
\usepackage[unicode,colorlinks=true,linkcolor=teal!55!black,urlcolor=teal!55!black]{hyperref}
\hypersetup{pdftitle={CLCNet 实现详解},pdfsubject={相机与激光雷达外参的置信度门控迭代}}
\definecolor{ink}{HTML}{1B343D}
\definecolor{accent}{HTML}{006B69}
\definecolor{pale}{HTML}{EBF5F2}
\definecolor{muted}{HTML}{47616A}
\definecolor{warm}{HTML}{AA5326}
\urlstyle{tt}
\newcommand{\code}[1]{\allowbreak\nolinkurl{#1}\allowbreak}
\newcommand{\topic}[1]{\par\vspace{.7em}{\large\sffamily\bfseries\color{accent}#1}\par\nobreak\vspace{.1em}}
\newcommand{\note}[1]{\par\smallskip\noindent\colorbox{pale}{\parbox{\dimexpr\linewidth-2\fboxsep}{\strut #1\strut}}\par}
\newcommand{\src}[1]{\par\vfill{\raggedright\small\color{muted}源码定位\quad #1\par}}
\newcommand{\chapterpage}[2]{\clearpage\section{#1}\noindent{\color{muted}#2}\par\medskip}
\newcommand{\Gt}{G}
\newcommand{\SE}{\mathrm{SE}(3)}
\newcolumntype{Y}{>{\raggedright\arraybackslash}X}
\setlength{\parindent}{0pt}
\setlength{\parskip}{.36em}
\setlength{\tabcolsep}{5pt}
\renewcommand{\arraystretch}{1.24}
\ctexset{section={format=\Large\sffamily\bfseries\color{ink},name={},number=\arabic{section},aftername=\quad,beforeskip=0pt,afterskip=.55em}}
\pagestyle{fancy}\fancyhf{}
\fancyhead[L]{\small\sffamily\color{muted}CLCNet 实现详解}
\fancyhead[R]{\small\color{muted}相机与激光雷达外参标定}
\fancyfoot[R]{\small\color{muted}\thepage}
\renewcommand{\headrulewidth}{.3pt}
\setlength{\fboxsep}{7pt}
\raggedbottom
\begin{document}
\pdfbookmark[1]{从一帧数据开始}{overview}
{\Huge\sffamily\bfseries\color{ink}CLCNet 实现详解}\par
{\Large\color{accent}从一帧图像和点云，到一次外参更新}\par\medskip
相机拍到一根灯杆，激光雷达也扫到了它。若按当前外参把激光点画回照片，点落在杆的右边，网络能利用这种错位估计一个三维修正。CLCNet 在 LCCNet 的平移和旋转输出旁加了一个分数，用它决定这一步修正该走多远，或者该不该走。灯杆只是例子，代码从整幅特征图学习匹配，无需检测灯杆或手工标注对应点。

\topic{先看推理时的一圈}
\begin{center}
\begin{tikzpicture}[font=\small,>=Stealth,box/.style={draw=accent,rounded corners=2pt,fill=pale,align=center,minimum height=11mm,text width=30mm},link/.style={->,thick,color=accent}]
\node[box] (rgb) at (0,0.76) {同一帧 RGB};
\node[box] (dep) at (0,-0.76) {当前深度图};
\node[box] (net) at (4.1,0) {特征匹配与回归};
\node[box] (gate) at (8.4,0) {整步修正与分数\\$\Delta t,\ \Delta q,\ c$};
\draw[link] (rgb.east) -- (net.west);
\draw[link] (dep.east) -- (net.west);
\draw[link] (net) -- (gate);
\draw[link] (gate.south) -- ++(0,-1.42) -| (dep.south);
\node[fill=white,inner sep=2pt,text=accent] at (4.1,-1.98) {接受并还有下一轮：用累计修正重投影点云};
\end{tikzpicture}
\end{center}
一旦拒绝，当前样本保留上一次接受的结果，结束本帧修正。整圈里 RGB 固定；点在图像中的落点随累计修正变化，所以第二轮看到的是一张新的深度图。

\topic{读代码时始终分清三个量}
\begin{tabularx}{\linewidth}{@{}lY@{}}\toprule
符号 & 含义 \\\midrule
$T_{CL}$ & 激光雷达系到相机系的外参。它回答“同一物理点在相机坐标系的坐标是多少”。\\
$\Delta_k$ & 第 $k$ 次前向给出的候选刚体变换，包括三维平移和三维旋转。\\
$A_k$ & 第 $k$ 步后的累计估计。每个样本从单位矩阵起步，代码用它的逆修正错位点云。\\\bottomrule
\end{tabularx}
\topic{阅读顺序}
第 \pageref{geo} 页用一个点算出像素错位；第 \pageref{net} 页看网络怎样找偏移；第 \pageref{loss}--\pageref{train} 页讲分数和训练；第 \pageref{eval} 页手算两步推理。只想先抓住工作原理，可以先读本页和第 \pageref{eval} 页，再回头看公式。
\src{主入口为 \code{CLCNet/tools/train.py} 与 \code{CLCNet/tools/evaluate.py}。下文默认列向量、平移单位米，示例数字均单独标明。}

\chapterpage{数据与坐标：错位的深度图从哪里来}{问题：只有一套已标定的数据，怎样得到既有错位、又知道正确答案的训练样本？}\label{geo}
\topic{先得到正确对齐的点，再人为制造可控误差}
数据集读取同帧 RGB、LiDAR 点云和标定文件。从投影矩阵 \code{P2} 取左侧 $3\times3$ 作为内参 $K$，从 \code{Tr_velo_to_cam} 取外参 $T_{CL}$。当前 VoD 加载器读取本划分第一帧的标定文件并在该划分复用。把原始点写作齐次列向量 $p_L=(x,y,z,1)^\top$，正确对齐的相机系点是 $p_C=T_{CL}p_L$。

按当前量程分别抽取三个平移分量及三个 XYZ 欧拉角，组成随机刚体变换 $H$。代码将其求逆得到监督目标 $\Gt=H^{-1}$，再对已经对齐的点施加 $\Gt$。记所有正确点组成的矩阵为 $P_C$，则
\[
 P_0=\Gt P_C,\qquad P(A)=A^{-1}P_0,\qquad P(\Gt)=P_C.
\]
第一式产生错位云，第二式是重投影前的点云，第三式提供了核对方向的办法。代码中的 \code{rotate_back(P,M)} 做 $MP$，\code{rotate_forward(P,M)} 做 $M^{-1}P$。只看函数名容易把方向颠倒。

\topic{把错位转换为网络看得见的像素偏移}
对相机系点 $(X,Y,Z)$，针孔投影给出
\[
 u=f_x X/Z+c_x,\qquad v=f_y Y/Z+c_y .
\]
图像先缩放到 $384\times1280$，再到网络的 $256\times512$（高$\times$宽），相应缩放焦距和主点。深度直接用缩放后的内参在网络分辨率上生成。保留视野内且 $0<Z<80\,\mathrm m$ 的点（默认上限），将 $Z/80$ 写入整数像素，空像素为零。多个点命中同一像素时，当前实现按索引赋值，没有显式选择最近点。

\topic{用一个点核对方向}
设网络分辨率下 $f_x=500$ 像素、$c_x=256$ 像素，正确点是 $(0,0,10)\,\mathrm m$。若 $\Gt$ 仅沿相机 $x$ 轴平移 $+0.20\,\mathrm m$，错位点为 $(0.20,0,10)$，其水平位置从 $256$ 变为 $266$。网络若精确输出 $A=\Gt$，重投影使用 $A^{-1}$ 把点移回 $x=0$，像素回到 $256$。这个例子也说明，同样的平移对近处点造成的像素偏移更大。
\note{若实际输入点已由当前外参 $T_{\rm cur}$ 转到相机系，按这里的约定，等效修正后的外参为 $A^{-1}T_{\rm cur}$。现有评估脚本计算并评价 $A$，实际写回标定文件属于部署端操作。}
\src{\code{LCCNet/DatasetVoD.py}、\code{LCCNet/utils.py}、\code{CLCNet/clcnet/utils/geometry.py}，以及 \code{train.py} 的 \code{build_depth_batch}。}

\chapterpage{网络：从局部匹配到三个输出}{问题：两幅图的数值和外观完全不同，网络靠什么推断三维旋转和平移？}\label{net}
\topic{两条编码支路}
RGB 支路用带预训练权重的 ResNet-18 提取外观特征；深度支路用单通道卷积和残差块提取稀疏深度的形状特征。默认 \code{use_feat_from=1}，实际进入匹配层的是两路最深的 $1/32$ 分辨率特征。把批量维度省略后，张量按“通道$\times$高$\times$宽”排列：
\begin{center}\begin{tabular}{@{}lll@{}}\toprule
对象 & 形状 & 内容 \\\midrule
RGB / 投影深度 & $3\times256\times512$ / $1\times256\times512$ & 同帧两种观测 \\
RGB / 深度特征 & 各 $512\times8\times16$ & 各网格位置的描述向量 \\
局部代价体积 & $81\times8\times16$ & 各候选位移的相关分数 \\
密集拼接结果 & $529\times8\times16$ & 匹配分数与卷积特征 \\
共享表示 & $512$ & 三个输出头的共同输入 \\\bottomrule
\end{tabular}\end{center}
\topic{为什么有 81 个通道}
对 RGB 特征网格的一个位置 $s$，在深度特征图上分别试探水平和竖直方向 $-4$ 到 $+4$ 的偏移 $\delta$。每个偏移都计算一组向量的平均内积：
\[
 V(\delta,s)=\frac{1}{512}\sum_{j=1}^{512}
 F_{\mathrm{rgb},j}(s)\,F_{\mathrm{depth},j}(s+\delta),\qquad \delta\in\{-4,\ldots,4\}^2 .
\]
这给出 $9\times9=81$ 个候选。偏移单位是特征网格格数。网络保留全部分数，让后续卷积结合多个位置的匹配分布推断整帧位姿。五个卷积块各追加 $128,128,96,64,32$ 个通道，总数为 $529$；展平后由全连接层生成共享向量。

\topic{CLCNet 增加的分支接在哪里}
LCCNet 从共享向量分别回归 $\Delta t\in\mathbb R^3$ 和四元数 $\Delta q=(w,x,y,z)$，并把四元数归一化。CLCNet 复用同一层激活后的特征，新增 $512\to256\to1$ 分支，最后用 Sigmoid 得到 $c\in(0,1)$。每个样本、每轮得到一个 $c$，用于整个位姿修正。初始化时末层权重为零、偏置为 1，因此新头最初总给出 $\sigma(1)\approx0.731$；只有训练后，分数才会随输入变化。
\src{\code{LCCNet/models/LCCNet.py} 的 \code{forward}；\code{LCCNet/models/correlation_package/correlation.py}；\code{CLCNet/clcnet/models/network.py} 的 \code{ConfLCCNet}。}

\chapterpage{监督：一个置信分数是怎样学出来的}{问题：没有人工给每帧打“可信”或“不可信”的标签，置信度头如何训练？}\label{loss}
\topic{用已知位姿误差构造软标签}
记送进损失函数的估计为 $\widehat A=(\hat t,\hat q)$，目标为 $\Gt=(t^*,q^*)$。令 $d_t=\|\hat t-t^*\|_2$，令 $d_r$ 为预测和目标间的最小相对转角。训练入口将当前档的平移、旋转扰动上限用作尺度 $s_t,s_r$。若 $d_t,s_t$ 均用米，$d_r^\circ,s_r$ 均用度，则
\[
 c^*=\exp\!\left[-\frac12\!\left(\frac{d_t}{s_t}+\frac{d_r^\circ}{s_r}\right)\right].
\]
例如第 4 档 $s_t=0.2\,\mathrm m,\ s_r=2^\circ$，若 $d_t=0.10\,\mathrm m,\ d_r^\circ=1^\circ$，则 $c^*=e^{-0.5}\approx0.607$。误差越小，标签越接近 1。标签由当前预测计算并停止梯度，置信头用二元交叉熵 $\mathrm{BCE}(c,c^*)$ 拟合这个连续值。推理时仅使用网络给出的 $c$。

\topic{总损失中的每一项在管什么}
以下按训练脚本的默认权重写，所有项最终取批均值。$d_r^{\rm rad}$ 表示弧度制转角。
\begin{align*}
 L_{\rm pose}&=2\sum_{j=1}^{3}\operatorname{SmoothL1}(\hat t_j-t_j^*)+d_r^{\rm rad},\\
 L_{\rm pc}&=\frac1N\sum_{i=1}^N
 \big\|(\Gt^{-1}\widehat A)^{-1}p_i-p_i\big\|_2,\quad p_i\in P_C,\\
 L_{\rm gate}&=(1-c)\big(\|\hat t\|_2/s_t+\angle_{\rm rad}(\hat q,q_I)\big),\\
 L&=0.5 L_{\rm pose}+0.5 L_{\rm pc}+0.2\,\mathrm{BCE}(c,c^*)+0.05 L_{\rm gate}.
\end{align*}
位姿项对齐平移和旋转。点云项用同一批点比较错误变换引起的三维位移；$\widehat A=\Gt$ 时每个点都回到原处。按点平均再按样本平均，有助于把点数不同的帧放在一起训练。质量配方至多抽取 16\,384 个点来计算此项。

门控惩罚把低分与大幅修正相乘。同样的位姿幅值下，提高 $c$ 也能降低这一项，因此要和置信监督一起看。角度在这里用弧度；只有构造 $c^*$ 时才按角度制量程归一化。
\note{怎样解释 $c=0.8$？它是模型对上述误差软标签的估计。把它解释成“执行本步有 80\% 的概率改善外参”，还需要逐步改善标签及相应的概率验证。}
\src{\code{CLCNet/clcnet/losses/confidence.py} 的 \code{pose_error_target_conf} 和 \code{ConfGatedCombinedLoss}；角距离见 \code{LCCNet/quaternion_distances.py}。}

\chapterpage{训练：一次前向和两步展开分别学什么}{问题：同一个置信分数既进损失，又参与位姿更新，这两条路径怎样相接？}\label{train}
\topic{单步训练}
当 \code{train_iters} 为 1 时，输入 $(I,D_0)$，模型给出原始预测 $(\Delta t_1,\Delta q_1,c_1)$。损失函数直接比较原始位姿与 $\Gt$，并用原始位姿误差生成 $c_1^*$。此分支不会先把预测平移乘以 $c_1$。

\topic{多步训练中的软门控}
令单位四元数 $q_I=(1,0,0,0)$。第 $k$ 步对平移做缩放，对旋转做线性混合后归一化：
\[
 \widetilde t_k=c_k\Delta t_k,\qquad
 \widetilde q_k=\operatorname{normalize}\!\left((1-c_k)q_I+c_k\Delta q_k\right).
\]
比如 $\Delta t_x=0.16\,\mathrm m,\ c=0.8$，训练时实际计入的平移为 $0.128\,\mathrm m$。旋转采用四元数归一化插值，通常不能直接把旋转角乘以 $c$。另外，$q$ 和 $-q$ 表示同一旋转，但代码插值前没有统一四元数符号，两种表示的中间值可能不同。

\topic{把默认两步训练完整展开}
\begin{align*}
 (t_1,q_1,c_1)&=f_\theta(I,D_0),&
 A_1&=\widetilde\Delta_1,\\
 D_1&=\Pi_K(A_1^{-1}P_0),&
 (t_2,q_2,c_2)&=f_\theta(I,D_1),\\
 A_2&=\widetilde\Delta_2 A_1,&
 L&=L(P_C,\Gt,A_2,c_2).
\end{align*}
这里 $\widetilde\Delta_k$ 由 $\widetilde t_k,\widetilde q_k$ 组成，$\Pi_K$ 包括投影和深度归一化。两次调用共享参数 $\theta$。每次用原始错位云 $P_0$ 和累计矩阵重建深度。代码左乘累计矩阵；带旋转时不能随意交换乘法顺序。深度写入整数像素，横纵像素索引的变化不提供连续的图像采样梯度。

\note{实际传入损失的是最终累计位姿 $A_2$ 和末步分数 $c_2$。$c_2^*$ 衡量两步累计结果离真值有多远。第一步的 $c_1$ 可以通过软门控影响累计位姿损失；主训练循环没有为每一步各算一次置信度 BCE。}

例如两次尝试都结束在 $10$ cm，且旋转误差相同。若一次从 $30$ cm 改到 $10$ cm，另一次从 $2$ cm 变为 $10$ cm，所得末端误差标签相同，更新带来的效果却相反。这是理解当前门控能力边界的关键。
\src{\code{CLCNet/tools/train.py} 的 \code{train_one_epoch}；软门控与位姿组装见 \code{CLCNet/clcnet/models/gated_refine.py} 的 \code{soft_gate_pose}、\code{compose_delta}。}

\chapterpage{推理：手算一次接受和一次拒绝}{问题：运行时没有真值，系统怎样决定继续修正，最后又返回什么？}\label{eval}
\topic{一帧内部，最多做三次尝试}
独立评估入口默认最多迭代 3 步、阈值 $\tau=0.5$。起初 $A_0=I_4$，样本处于活动状态。每轮网络从 RGB 和当前深度给出未经缩放的候选 $\Delta_k$ 与分数 $c_k$：
\begin{itemize}
\item 若样本仍活动且 $c_k\ge\tau$，接受整步，令 $A_k=\Delta_k A_{k-1}$。若还有下一轮，重新投影 $A_k^{-1}P_0$ 生成深度图。
\item 若 $c_k<\tau$，保留已有 $A_{k-1}$，该样本此后保持冻结。若第一步就被拒绝，最终累计估计是单位矩阵，外参保持原值。
\end{itemize}
实现用批内布尔掩码控制接受。其他样本仍活动时，已冻结样本也可能参与批量前向，只是它的估计不再变化。\code{steps} 统计真正接受的步数；它与前向调用次数、端到端延迟是三个不同指标。

\topic{沿用第 \pageref{geo} 页的灯杆，观察一个像素}
下面是假设的网络输出，仅为手算示例。真实错位为相机 $x$ 轴 $+0.20\,\mathrm m$，初始像素 $u=266$，正确像素 $u^*=256$；各步旋转均为单位旋转。
\begin{center}\begin{tabular}{@{}lrrrrl@{}}\toprule
时刻 & 候选 $\Delta t_x$ & $c$ & 累计 $A_x$ & 像素 $u$ & 动作\\\midrule
初始 & --- & --- & $0.00$ m & $266$ & 等待预测\\
第 1 步 & $+0.16$ m & $0.80$ & $+0.16$ m & $258$ & 接受全量\\
第 2 步 & $+0.10$ m & $0.35$ & $+0.16$ m & $258$ & 拒绝并冻结\\\bottomrule
\end{tabular}\end{center}
第一步后剩余 $0.20-0.16=0.04\,\mathrm m$，所以 $u=256+500\times0.04/10=258$。第二步若强行执行，点将落在 $u=253$，离正确像素从 2 像素变成 3 像素。在这个构造例中，拒绝保留了较好的结果，最终 \code{steps=1}。
\begin{center}
\begin{tikzpicture}[x=.72cm,y=1cm,>=Stealth,font=\small]
\draw[->,color=muted] (-.6,0)--(14.4,0) node[right] {$u$};
\foreach \x/\u in {0/253,3/256,5/258,13/266}{\draw (\x,.07)--(\x,-.07) node[below]{\u};}
\draw[accent,very thick,->] (13,.45)--(5,.45) node[midway,above]{第 1 步接受};
\draw[warm,dashed,thick,->] (5,.45)--(0,.45) node[midway,above]{第 2 步若执行};
\draw[accent,thick] (3,-.35)--(3,-.49) node[below]{灯杆真实位置};
\fill[accent] (5,0) circle(2.1pt);\fill[muted] (13,0) circle(2.1pt);
\end{tikzpicture}
\end{center}

\note{同样是 $c=0.80$ 和 $+0.16$ m，训练软门控走 $+0.128$ m，推理硬门控走 $+0.16$ m。一帧结束后，评估脚本将累计估计与该帧真值比较。下一帧重新初始化；持续维护车辆外参还需跨帧状态传递。}
\src{\code{CLCNet/tools/evaluate.py} 的 \code{active}、\code{gate}、\code{steps} 与 \code{acc}；重投影见 \code{project_batch_after_rt}。}

\chapterpage{读训练配置和日志：五档、两步和三步}{问题：文件夹叫 iter1 至 iter5，为什么日志里又能看到 train\_iters=2 和 steps=3.00？}\label{config}
\topic{五档描述训练样本的难度}
\begin{center}\begin{tabular}{@{}lrrr@{}}\toprule
训练档位 & 各轴平移采样范围 & 各欧拉角采样范围 & 质量配方 epochs\\\midrule
\code{iter1} & $\pm1.5$ m & $\pm20^\circ$ & 50\\
\code{iter2} & $\pm1.0$ m & $\pm10^\circ$ & 40\\
\code{iter3} & $\pm0.5$ m & $\pm5^\circ$ & 40\\
\code{iter4} & $\pm0.2$ m & $\pm2^\circ$ & 40\\
\code{iter5} & $\pm0.1$ m & $\pm1^\circ$ & 50\\\bottomrule
\end{tabular}\end{center}
表中是 \code{train_5stage.py --quality} 的代码配置。第一档可以加载已有 LCCNet 权重，后续各档以前一档最佳 CLCNet 权重初始化。表中范围是对采样前的平移分量、欧拉角分量分别限幅，三维平移范数和合成转角可以更大。训练时随机采样；验证集用种子 123 固定每帧扰动，以便同条件比较。

\topic{两步和三步都发生在同一帧}
\code{train_iters=2} 指当前档训练时用同一组权重展开两轮软门控。\code{eval_iters=3} 指评估时给同一帧最多三次硬门控尝试。独立评估脚本只接收一个 \code{--weights}；每轮重复调用所加载的网络。假设传入 \code{iter1} 权重跑三轮，这三轮全部使用第一档模型。

\topic{把一条真实日志翻译成一句话}
第一档续训日志中有一次验证记录：平移中位误差 $21.794$ cm、旋转中位误差 $2.309^\circ$、平均接受步数 $3.00$。平移误差是每帧 $\|\hat t-t^*\|_2$ 转为厘米后取中位数；旋转是相对转角转为度后取中位数。平均步数显示到两位小数，说明该次验证几乎用满了三步。单凭这条记录，无法确定哪些候选更新有害，也无法量出门控挽回了多少误差。

工作区另有 $5.24$ cm / $0.438^\circ$ 与 $7.71$ cm / $0.924^\circ$ 两组数据；基线交底书将它们分别记为 LCCNet 最细单档与五档级联复现成绩。和上面的 CLCNet 第一档日志比较时，需先对齐权重、扰动量程及迭代协议。
\src{\code{CLCNet/tools/train_5stage.py} 的 \code{STAGES} 与 \code{QUALITY_EPOCHS}；\code{CLCNet/outputs/logs/vod_quality/train_iter1_resume.log} 第 29 行；《LCCNet相机激光雷达外参标定交底书》第 103 行。}

\chapterpage{回到源码：顺着一帧数据读}{实际运行以工具脚本的调用链为准。下面的入口按数据流排列；路径从工作区根目录算起。}\label{map}
\begin{tabularx}{\linewidth}{@{}p{28mm}Y@{}}\toprule
想弄懂什么 & 打开哪里，追踪哪个量\\\midrule
真值怎么来 & \code{LCCNet/DatasetVoD.py} 的 \code{__getitem__}：跟踪 \code{pc_cam}、\code{invert_pose}、\code{tr_error}、\code{rot_error}。\\
深度怎么造 & \code{CLCNet/tools/train.py} 的 \code{build_depth_batch}，接着看 \code{clcnet/utils/geometry.py} 的投影与重投影（后一路径相对 \code{CLCNet/}）。\\
两路如何匹配 & \code{LCCNet/models/LCCNet.py} 的 \code{forward}：从 \code{c16}、\code{c26} 追到 \code{corr6}，再看 \code{fc1} 和两个位姿头。\\
分数从哪里出 & \code{CLCNet/clcnet/models/network.py} 的 \code{ConfLCCNet.forward}：跟踪共享特征 \code{x} 到 \code{conf}。\\
每步怎样累计 & \code{CLCNet/tools/train.py} 的 \code{train_one_epoch}：顺着 \code{soft_gate_pose}、\code{acc}、\code{conf_last} 追到 \code{loss_fn}。\\
分数怎样监督 & \code{CLCNet/clcnet/losses/confidence.py}：看 \code{conf_tgt}、\code{conf_loss}、\code{gate_pen} 分别怎样参与总损失。\\
何时拒绝 & \code{CLCNet/tools/evaluate.py} 的 \code{main}：看 \code{gate} 怎样控制 \code{delta_m}，以及 \code{active} 怎样冻结样本。\\\bottomrule
\end{tabularx}
\topic{三个可以检验理解的问题}
\textbf{1. 为什么第一次修正后还要重新投影？}\quad RGB 描述的是同一个场景，改变的是点云在图像中的位置。用新深度图，第二轮才看得到第一轮留下的错位。

\textbf{2. 同一步输出四个旋转数，为什么总共还是六个自由度？}\quad 单位四元数的四个分量受单位范数约束，只表示三个旋转自由度；加上三维平移得到六个。置信分数属于更新控制量。

\textbf{3. 怎样判断门控确实帮了忙？}\quad 在同一组权重和固定扰动上，对每个候选分别记录执行前后的平移误差、旋转误差和分数，再比较强制接受与按阈值接受后的最终误差。这个比较能把回归器本身的效果和拒绝动作的贡献分开。
\src{本文按工作区源码及已有日志编写，配置表描述脚本计划，数值示例用于解释计算过程。}
\end{document}
